% Auto-generated from raw.txt
% Usage:
%   \vocabentry{word}{/ipa/ (pos)}{definition}{example}{note}

\vocabentry{Adolescence}
{/ˌædəˈlesns/ (n)}
{The period following the onset of puberty during which a young person develops from a child into an adult.}
{Adolescence is a turbulent time of life for many teenagers.}
{Giai đoạn thanh thiếu niên từ trẻ con sang người lớn.}

\vocabentry{Puberty}
{/ˈpjuːbərti/ (n)}
{The period during which adolescents reach sexual maturity and become capable of reproduction.}
{Most girls reach puberty between the ages of 10 and 14.}
{Giai đoạn dậy thì của trẻ vị thành niên.}

\vocabentry{Generation Gap}
{/ˌdʒenəˈreɪʃn ɡæp/ (n)}
{A difference of opinions between one generation and another regarding beliefs, politics, or values.}
{The generation gap often leads to misunderstandings between parents and children.}
{Khoảng cách thế hệ về quan điểm và tư duy.}

\vocabentry{Peer pressure}
{/ˈpɪər preʃər/ (n)}
{Influence from members of one's peer group to behave in a manner similar to them.}
{Many teenagers start smoking due to peer pressure.}
{Áp lực phải làm theo hoặc giống với bạn bè đồng trang lứa.}

\vocabentry{Juvenile}
{/ˈdʒuːvənl/ (adj)}
{For or relating to young people who are not yet adults.}
{The town has seen a rise in juvenile crime over the past year.}
{Liên quan đến người vị thành niên, thường dùng trong ngữ cảnh pháp luật.}

\vocabentry{Delinquency}
{/dɪˈlɪŋkwənsi/ (n)}
{Minor crime, especially that committed by young people.}
{Poverty is often linked to high rates of juvenile delinquency.}
{Sự phạm pháp hoặc hành vi sai trái của người trẻ tuổi.}

\vocabentry{Socialization}
{/ˌsoʊʃələˈzeɪʃn/ (n)}
{The process of learning to behave in a way that is acceptable to society.}
{Schools play a vital role in the socialization of children.}
{Quá trình hòa nhập và học hỏi các quy tắc xã hội.}

\vocabentry{Alienation}
{/ˌeɪliəˈneɪʃn/ (n)}
{The state or experience of being isolated from a group or an activity to which one should belong.}
{Unemployment can lead to a sense of social alienation among young people.}
{Cảm giác bị cô lập hoặc xa lánh khỏi cộng đồng.}

\vocabentry{Rebellion}
{/rɪˈbeljən/ (n)}
{An act of open resistance to an established government, ruler, or parental authority.}
{Teenagers often go through a period of rebellion against their parents' rules.}
{Sự nổi loạn hoặc chống đối lại các quy tắc/cha mẹ.}

\vocabentry{Conformity}
{/kənˈfɔːrməti/ (n)}
{Behavior in accordance with socially accepted standards or conventions.}
{There is a lot of pressure for conformity in high school.}
{Sự tuân thủ theo các chuẩn mực và kỳ vọng của xã hội.}

\vocabentry{Identity crisis}
{/aɪˈdentəti ˌkraɪsɪs/ (n)}
{A period of uncertainty and confusion in which a person's sense of identity becomes insecure.}
{Many adolescents experience an identity crisis as they try to find their place in the world.}
{Khủng hoảng danh tính, không rõ mình là ai.}

\vocabentry{Self-esteem}
{/ˌself ɪˈstiːm/ (n)}
{Confidence in one's own worth or abilities; self-respect.}
{Low self-esteem can prevent young people from pursuing their dreams.}
{Lòng tự trọng hoặc sự tự tin vào giá trị của bản thân.}

\vocabentry{Narcissism}
{/ˈnɑːrsɪsɪzəm/ (n)}
{Excessive interest in or admiration of oneself and one's physical appearance.}
{Some critics argue that social media encourages narcissism in Gen Z.}
{Sự tự luyến, quá yêu bản thân và khao khát sự chú ý.}

\vocabentry{Millennials}
{/mɪˈleniəlz/ (n)}
{People reaching young adulthood in the early 21st century, typically born between 1981 and 1996.}
{Millennials are known for being the first generation to grow up with the internet.}
{Thế hệ Millennials, hay thế hệ Y.}

\vocabentry{Gen Z}
{/ˌdʒen ˈziː/ (n)}
{The generation born between the late 1990s and the early 2010s.}
{Gen Z is often characterized by its digital fluency and social activism.}
{Thế hệ Z, lớn lên trong kỷ nguyên số toàn diện.}

\vocabentry{Digital natives}
{/ˌdɪdʒɪtl ˈneɪtɪvz/ (n)}
{People born or brought up during the age of digital technology and therefore familiar with it from early childhood.}
{Unlike their parents, Gen Z are true digital natives.}
{Người bản địa số.}

\vocabentry{Cyberbullying}
{/ˈsaɪbərbʊliɪŋ/ (n)}
{The use of electronic communication to bully a person, typically by sending messages of an intimidating nature.}
{Schools are taking measures to educate students about the dangers of cyberbullying.}
{Bắt nạt qua mạng.}

\vocabentry{Screen addiction}
{/ˈskriːn əˌdɪkʃn/ (n)}
{The compulsive use of digital devices such as smartphones or computers.}
{Screen addiction can lead to sleep deprivation and poor academic performance.}
{Chứng nghiện sử dụng điện thoại và các thiết bị điện tử.}

\vocabentry{Career orientation}
{/kəˈrɪr ˌɔːriənˈteɪʃn/ (n)}
{The process of assisting young people in choosing and preparing for a profession.}
{Good career orientation in high school helps students make better college choices.}
{Hoạt động định hướng nghề nghiệp.}

\vocabentry{Job insecurity}
{/ˌdʒɑːb ˌɪnsɪˈkjʊrəti/ (n)}
{The state of knowing that your job is not permanent or may be lost at any time.}
{Many young graduates face job insecurity in today's competitive market.}
{Sự bấp bênh và thiếu ổn định trong công việc.}

\vocabentry{Underemployment}
{/ˌʌndərɪmˈplɔɪmənt/ (n)}
{The condition of being employed at a job that is inadequate to one's training.}
{Underemployment is a serious issue for many highly educated young people.}
{Tình trạng làm việc dưới năng lực.}

\vocabentry{Gig economy}
{/ˈɡɪɡ ɪˌkɑːnəmi/ (n)}
{A labor market characterized by short-term contracts or freelance work.}
{The gig economy offers flexibility but lacks the benefits of a stable job.}
{Nền kinh tế việc làm tự do hoặc thời vụ.}

\vocabentry{Student debt}
{/ˈstjuːdnt det/ (n)}
{Financial obligations accumulated by students to cover the costs of their education.}
{High levels of student debt are delaying major life milestones like buying a home.}
{Các khoản nợ vay để học đại học của sinh viên.}

\vocabentry{Academic pressure}
{/ˌækəˈdemɪk ˈpreʃər/ (n)}
{Stress and anxiety resulting from the high demands of school performance.}
{Excessive academic pressure can lead to mental health problems in teenagers.}
{Áp lực học hành thi cử.}

\vocabentry{Extracurricular}
{/ˌekstrəkəˈrɪkjələr/ (adj)}
{An activity at a school or college pursued in addition to the normal course of study.}
{Universities look for students who are active in extracurricular activities.}
{Hoạt động ngoại khóa.}

\vocabentry{Gap year}
{/ˈɡæp jɪər/ (n)}
{A year taken by a student as a break between school and university.}
{He decided to take a gap year to travel and volunteer in Africa.}
{Năm nghỉ để trải nghiệm trước khi vào đại học.}

\vocabentry{Volunteering}
{/ˌvɑːlənˈtɪrɪŋ/ (n)}
{Freely offering to do something or help others, often for a community or organization.}
{Volunteering at a local hospital provides valuable experience for medical students.}
{Hoạt động tình nguyện.}

\vocabentry{Soft skills}
{/ˈsɔːft skɪlz/ (n)}
{Personal attributes that enable someone to interact effectively with others.}
{Employers value soft skills just as much as technical qualifications.}
{Kỹ năng mềm như giao tiếp, teamwork.}

\vocabentry{Interpersonal}
{/ˌɪntərˈpɜːrsənl/ (adj)}
{Relating to relationships or communication between people.}
{He has excellent interpersonal skills and gets along with everyone.}
{Liên quan đến giao tiếp cá nhân.}

\vocabentry{Critical thinking}
{/ˌkrɪtɪkl ˈθɪŋkɪŋ/ (n)}
{The objective analysis of an issue in order to form a judgment.}
{The school curriculum aims to develop students' critical thinking abilities.}
{Tư duy phản biện.}

\vocabentry{Empathy}
{/ˈempəθi/ (n)}
{The ability to understand and share the feelings of another person.}
{Reading literature helps young people develop empathy for others.}
{Sự thấu cảm.}

\vocabentry{Mindfulness}
{/ˈmaɪndflnəs/ (n)}
{A mental state achieved by focusing one's awareness on the present moment.}
{Practicing mindfulness can help students manage exam-related stress.}
{Sự chánh niệm.}

\vocabentry{Mental health}
{/ˌmentl ˈhelθ/ (n)}
{A person's condition with regard to their psychological and emotional well-being.}
{There is a growing awareness of the importance of mental health among youth.}
{Sức khỏe tâm thần.}

\vocabentry{Burnout}
{/ˈbɜːrnaʊt/ (n)}
{State of emotional, physical, and mental exhaustion caused by excessive stress.}
{Studying for 12 hours a day eventually led to total burnout.}
{Sự kiệt sức.}

\vocabentry{Depression}
{/dɪˈpreʃn/ (n)}
{A serious medical condition in which a person feels very sad and hopeless.}
{It is important to seek professional help if you are suffering from depression.}
{Chứng bệnh trầm cảm.}

\vocabentry{Anxiety}
{/æŋˈzaɪəti/ (n)}
{A feeling of worry, nervousness, or unease about an uncertain outcome.}
{Social anxiety can make it difficult for young people to make friends.}
{Sự lo âu.}

\vocabentry{Loneliness}
{/ˈloʊnlinəs/ (n)}
{The sadness that is caused by having no friends or being by oneself.}
{Despite having many followers online, she often felt a sense of loneliness.}
{Sự cô đơn.}

\vocabentry{Substance abuse}
{/ˈsʌbstəns əˌbjuːs/ (n)}
{Overindulgence in or dependence on an addictive substance.}
{Education is key to preventing substance abuse among teenagers.}
{Lạm dụng chất gây nghiện.}

\vocabentry{Alcoholism}
{/ˈælkəhɔːlɪzəm/ (n)}
{An addiction to the consumption of alcoholic liquor.}
{Binge drinking at parties can lead to long-term alcoholism.}
{Chứng nghiện rượu.}

\vocabentry{Vaping}
{/ˈveɪpɪŋ/ (n)}
{The action of inhaling and exhaling the vapor produced by an electronic cigarette.}
{Vaping has become a major health concern in middle and high schools.}
{Việc hút thuốc lá điện tử.}

\vocabentry{Obesity}
{/oʊˈbiːsəti/ (n)}
{The state of being overweight to a degree that is unhealthy.}
{Childhood obesity is often linked to a diet high in processed foods.}
{Bệnh béo phì.}

\vocabentry{Sedentary}
{/ˈsednteri/ (adj)}
{Characterized by much sitting and little physical exercise.}
{A sedentary lifestyle increases the risk of heart disease.}
{Thói quen ít vận động.}

\vocabentry{Nutrition}
{/nuˈtrɪʃn/ (n)}
{The process of obtaining the food necessary for health and growth.}
{Proper nutrition is essential for brain development in young children.}
{Chế độ dinh dưỡng.}

\vocabentry{Body shaming}
{/ˈbɑːdi ˌʃeɪmɪŋ/ (n)}
{The practice of making critical comments about a person's body size or shape.}
{Body shaming on social media can have a devastating impact on self-esteem.}
{Miệt thị ngoại hình.}

\vocabentry{Eating disorder}
{/ˈiːtɪŋ dɪsˌɔːrdər/ (n)}
{Any of a range of psychological disorders characterized by abnormal eating habits.}
{Anorexia is a severe eating disorder that requires medical treatment.}
{Rối loạn ăn uống.}

\vocabentry{Self-harm}
{/ˌself ˈhɑːrm/ (n)}
{Deliberate injury to one's own body as a manifestation of psychological distress.}
{Parents should watch for signs of self-harm in their children.}
{Hành vi tự làm hại bản thân.}

\vocabentry{Suicide}
{/ˈsuːɪsaɪd/ (n)}
{The act of intentionally taking one's own life.}
{We need better support systems to prevent youth suicide.}
{Sự tự sát.}

\vocabentry{Peer group}
{/ˈpɪər ɡruːp/ (n)}
{A social group whose members have similar interests and age.}
{For many adolescents, the peer group is more influential than the family.}
{Nhóm bạn bè cùng trang lứa.}

\vocabentry{Social status}
{/ˌsoʊʃl ˈsteɪtəs/ (n)}
{A person's standing or importance in relation to other people within a society.}
{In high school, wearing brand-name clothes is often a way to gain social status.}
{Địa vị xã hội.}

\vocabentry{Activism}
{/ˈæktɪvɪzəm/ (n)}
{Vigorous campaigning to bring about political or social change.}
{Young people are leading the way in global climate activism.}
{Chủ nghĩa hoạt động/tích cực.}

\vocabentry{Climate anxiety}
{/ˈklaɪmət æŋˌzaɪəti/ (n)}
{Worry about the future of the environment due to climate change.}
{Many young people feel climate anxiety when they think about the future of the planet.}
{Nỗi lo lắng về biến đổi khí hậu.}

\vocabentry{Environmentalism}
{/ɪnˌ|vīrənˈmentlˌizəm/ (n)}
{Concern about and action aimed at protecting the environment.}
{Her environmentalism led her to start a recycling program at her school.}
{Chủ nghĩa bảo vệ môi trường.}

\vocabentry{Sustainability}
{/səˌsteɪnəˈbɪləti/ (n)}
{The avoidance of the depletion of natural resources in order to maintain balance.}
{Sustainability is a core value for many Gen Z consumers.}
{Tính bền vững.}

\vocabentry{Consumerism}
{/kənˈsuːmərɪzəm/ (n)}
{The preoccupation of society with the acquisition of consumer goods.}
{Fast fashion is a prominent example of modern consumerism.}
{Chủ nghĩa tiêu dùng.}

\vocabentry{Materialism}
{/məˈtɪriəlɪzəm/ (n)}
{Tendency to consider material possessions as more important than spiritual values.}
{Critics argue that advertising fuels materialism in young children.}
{Lối sống vật chất.}

\vocabentry{Influencer}
{/ˈɪnfluənsər/ (n)}
{A person who promotes products or lifestyles on social media to a large audience.}
{Many teenagers aspire to become a social media influencer as a career.}
{Người có sức ảnh hưởng.}

\vocabentry{Followers}
{/ˈfɑːloʊərz/ (n)}
{People who subscribe to the updates of another person on social media.}
{She has over a million followers on her travel blog.}
{Người theo dõi.}

\vocabentry{Engagement}
{/ɪnˈɡeɪdʒmənt/ (n)}
{Interaction (likes, shares, comments) content receives from its audience.}
{Companies look for high engagement rates when choosing influencers to work with.}
{Sự tương tác.}

\vocabentry{Trend}
{/trend/ (n)}
{A general direction in which something is developing or changing.}
{Keeping up with the latest fashion trends can be expensive.}
{Xu hướng.}

\vocabentry{Viral}
{/ˈvaɪrəl/ (adj)}
{Content that circulates rapidly and widely from one user to another.}
{The video of the dancing cat went viral within hours.}
{Lan truyền nhanh chóng.}

\vocabentry{Misinformation}
{/ˌmɪsɪnfərˈmeɪʃn/ (n)}
{False or inaccurate information that is spread, regardless of intent.}
{Social media companies are struggling to stop the spread of misinformation.}
{Thông tin sai lệch.}

\vocabentry{Privacy}
{/ˈpraɪvəsi/ (n)}
{The state or condition of being free from public attention.}
{You should be careful about your online privacy settings.}
{Sự riêng tư.}

\vocabentry{Surveillance}
{/sərˈveɪləns/ (n)}
{Close observation, especially related to digital tracking.}
{Many people are concerned about government surveillance of the internet.}
{Sự giám sát.}

\vocabentry{Artificial Intelligence}
{/ˌɑːrtɪfɪʃl ɪnˈtelɪdʒəns/ (n)}
{Computer systems able to perform tasks that normally require human intelligence.}
{AI is changing the way students research and write their essays.}
{Trí tuệ nhân tạo.}

\vocabentry{Automation}
{/ˌɔːtəˈmeɪʃn/ (n)}
{The use of machines or computers to perform tasks previously done by people.}
{Automation is a major concern for young people entering the manufacturing sector.}
{Sự tự động hóa.}

\vocabentry{Disruption}
{/dɪsˈrʌpʃn/ (n)}
{Radical change caused by technological innovation.}
{The rise of streaming services caused a massive disruption in the film industry.}
{Sự gián đoạn/xáo trộn đột ngột.}

\vocabentry{Innovation}
{/ˌɪnəˈveɪʃn/ (n)}
{A new method, idea, or product.}
{Innovation is essential for economic growth and solving global problems.}
{Sự đổi mới/sáng tạo.}

\vocabentry{Entrepreneurship}
{/ˌɑːntrəprəˈnɜːrʃɪp/ (n)}
{The activity of setting up a business and taking financial risks.}
{The university offers courses to encourage student entrepreneurship.}
{Khởi nghiệp.}

\vocabentry{Start-up}
{/ˈstɑːrt ʌp/ (n)}
{A newly established business with high growth potential.}
{Many successful start-ups were founded in university dorm rooms.}
{Công ty khởi nghiệp.}

\vocabentry{Global citizenship}
{/ˌɡloʊbl ˈsɪtɪzənʃɪp/ (n)}
{The idea that one's identity belongs to a worldwide community.}
{Global citizenship involves understanding and respecting diverse cultures.}
{Tinh thần công dân toàn cầu.}

\vocabentry{Intercultural}
{/ˌɪntərˈkʌltʃərəl/ (adj)}
{Relating to or involving different cultures.}
{Intercultural dialogue is necessary for global peace.}
{Tính chất liên văn hóa.}

\vocabentry{Tolerance}
{/ˈtɑːlərəns/ (n)}
{The ability to accept opinions or behavior that one does not agree with.}
{Schools should teach students the value of religious and cultural tolerance.}
{Sự khoan dung.}

\vocabentry{Diversity}
{/daɪˈvɜːrsəti/ (n)}
{The state of being diverse; variety in backgrounds.}
{We celebrate the cultural diversity of our local community.}
{Sự đa dạng.}

\vocabentry{Inclusion}
{/ɪnˈkluːʒn/ (n)}
{The practice of providing equal access to opportunities for all.}
{The company is committed to the inclusion of people with disabilities in its workforce.}
{Sự hòa nhập.}

\vocabentry{Discrimination}
{/dɪˌskrɪmɪˈneɪʃn/ (n)}
{Unjust treatment of different categories of people.}
{Discrimination on the basis of race is illegal in most countries.}
{Sự phân biệt đối xử.}

\vocabentry{Racism}
{/ˈreɪsɪzəm/ (n)}
{Prejudice or antagonism directed against someone of a different race.}
{Many young people are protesting against systemic racism.}
{Chủ nghĩa phân biệt chủng tộc.}

\vocabentry{Sexism}
{/ˈseksɪzəm/ (n)}
{Prejudice or discrimination on the basis of sex.}
{Sexism in the workplace can prevent women from reaching leadership positions.}
{Sự phân biệt giới tính.}

\vocabentry{Gender equality}
{/ˌdʒendər iˈkwɑːləti/ (n)}
{Access to rights or opportunities unaffected by gender.}
{Promoting gender equality is a key goal of the United Nations.}
{Bình đẳng giới.}

\vocabentry{Feminism}
{/ˈfemənɪzəm/ (n)}
{The advocacy of women's rights on the basis of equality of the sexes.}
{Feminism has evolved over the decades to address many different issues.}
{Chủ nghĩa nữ quyền.}

\vocabentry{Secularism}
{/ˈsekjələrɪzəm/ (n)}
{The principle of separation of the state from religious institutions.}
{Secularism is a fundamental principle of many modern democracies.}
{Chủ nghĩa thế tục.}

\vocabentry{Awareness}
{/əˈwernəs/ (n)}
{Knowledge or perception of a situation or fact.}
{The campaign aims to raise awareness about the importance of mental health.}
{Nhận thức.}

\vocabentry{Voting}
{/ˈvoʊtɪŋ/ (n)}
{Formal indication of a choice between candidates in an election.}
{Many organizations are working to increase youth voting rates.}
{Việc bỏ phiếu.}

\vocabentry{Protest}
{/ˈproʊtest/ (n)}
{A statement or action expressing disapproval of something.}
{Students organized a protest against the increase in tuition fees.}
{Cuộc biểu tình.}

\vocabentry{Democracy}
{/dɪˈmɑːkrəsi/ (n)}
{Government by the whole population, typically through representatives.}
{Freedom of speech is a cornerstone of any democracy.}
{Chế độ dân chủ.}

\vocabentry{Human rights}
{/ˌhjuːmən ˈraɪts/ (n)}
{Rights believed to belong justifiably to every person.}
{Everyone should have their basic human rights respected and protected.}
{Quyền con người.}

\vocabentry{Poverty}
{/ˈpɑːvərti/ (n)}
{The state of being extremely poor.}
{Poverty can prevent children from getting a good education.}
{Sự nghèo đói.}

\vocabentry{Inequality}
{/ˌɪnɪˈkwɑːləti/ (n)}
{Lack of equality in opportunities or wealth.}
{Income inequality is a major issue in many developed nations.}
{Sự bất bình đẳng.}

\vocabentry{Homelessness}
{/ˈhoʊmləsnəs/ (n)}
{The state of having no home.}
{Lack of affordable housing is a major cause of homelessness.}
{Tình trạng vô gia cư.}

\vocabentry{Unemployment}
{/ˌʌnɪmˈplɔɪmənt/ (n)}
{The state of being without a paid job.}
{High rates of youth unemployment can lead to social unrest.}
{Tình trạng thất nghiệp.}

\vocabentry{Inflation}
{/ɪnˈfleɪʃn/ (n)}
{A general increase in prices and fall in the purchasing value of money.}
{Inflation makes it difficult for young families to save money.}
{Lạm phát.}

\vocabentry{Cost of living}
{/ˌkɔːst əv ˈlɪvɪŋ/ (n)}
{The amount of money needed to sustain a certain level of living.}
{The high cost of living in the city center is forcing people to move further away.}
{Chi phí sinh hoạt.}

\vocabentry{Affordable}
{/əˈfɔːrdəbl/ (adj)}
{Inexpensive; reasonably priced.}
{We need more affordable housing options for young graduates.}
{Giá cả phải chăng.}

\vocabentry{Urbanization}
{/ˌɜːrbənaɪˈzeɪʃn/ (n)}
{The process of making an area more urban.}
{Rapid urbanization is putting pressure on city infrastructure.}
{Sự đô thị hóa.}

\vocabentry{Migration}
{/maɪˈɡreɪʃn/ (n)}
{Movement of people to a new area or country for work or life.}
{Migration from rural to urban areas is a global trend.}
{Sự di cư.}

\vocabentry{Globalization}
{/ˌɡloʊbələˈzeɪʃn/ (n)}
{Process by which businesses and organizations develop international influence.}
{Globalization has led to the spread of ideas and cultures across the world.}
{Sự toàn cầu hóa.}

\vocabentry{Cultural identity}
{/ˌkʌltʃərəl aɪˈdentəti/ (n)}
{The feeling of belonging to a group based on shared culture.}
{Maintaining your native language is key to preserving your cultural identity.}
{Bản sắc văn hóa.}

\vocabentry{Tradition}
{/trəˈdɪʃn/ (n)}
{Transmission of customs or beliefs from generation to generation.}
{It is a local tradition to celebrate the harvest festival with a large meal.}
{Truyền thống.}

\vocabentry{Heritage}
{/ˈherɪtɪdʒ/ (n)}
{Valued objects and qualities passed down from previous generations.}
{We must work to protect our national architectural heritage.}
{Di sản.}

\vocabentry{Customs}
{/ˈkʌstəmz/ (n)}
{A traditional way of behaving that is specific to a particular society.}
{Understanding local customs is important when traveling abroad.}
{Phong tục.}

\vocabentry{Folklore}
{/ˈfoʊklɔːr/ (n)}
{Traditional beliefs and stories passed on by word of mouth.}
{The book is based on Irish folklore and legends.}
{Văn hóa dân gian.}

\vocabentry{Superstition}
{/ˌsuːpərˈstɪʃn/ (n)}
{Credulous belief in supernatural beings or irrational luck.}
{Breaking a mirror is a common superstition that brings bad luck.}
{Sự mê tín dị đoan.}

\vocabentry{Values}
{/ˈvæljuːz/ (n)}
{Principles or standards of behavior; what is important in life.}
{Family values are very important in many Asian cultures.}
{Các giá trị.}

\vocabentry{Morality}
{/məˈræləti/ (n)}
{Principles concerning the distinction between right and wrong.}
{The film explores the complex morality of modern warfare.}
{Đạo đức cá nhân.}

\vocabentry{Ethics}
{/ˈeθɪks/ (n)}
{Moral principles that govern a person's behavior.}
{Business ethics require companies to be honest and transparent.}
{Luân lý/Đạo đức học.}

\vocabentry{Spirituality}
{/ˌspɪrɪtʃuˈæləti/ (n)}
{Quality of being concerned with the human spirit as opposed to material things.}
{Many people find peace through meditation and spirituality.}
{Sự tâm linh.}

\vocabentry{Religion}
{/rɪˈlɪdʒən/ (n)}
{Belief in and worship of a superhuman controlling power/God.}
{Religion plays a central role in the lives of many people.}
{Tôn giáo.}

\vocabentry{Faith}
{/feɪθ/ (n)}
{Complete trust or confidence in someone or something.}
{His faith helped him get through the difficult times in his life.}
{Đức tin/Niềm tin.}

\vocabentry{Community}
{/kəˈmjuːnəti/ (n)}
{A group of people living in the same place or sharing characteristics.}
{Volunteering is a great way to give back to your community.}
{Cộng đồng.}

\vocabentry{Dynamics}
{/daɪˈnæmɪks/ (n)}
{Forces which stimulate growth or change within a group.}
{The dynamics of the family changed after the birth of the new baby.}
{Động lực/Mối tương tác.}

\vocabentry{Parenting}
{/ˈperəntɪŋ/ (n)}
{The activity of bringing up a child as a parent.}
{Good parenting involves providing both love and discipline.}
{Việc nuôi dạy con cái.}

\vocabentry{Authoritarian}
{/əˌθɔːrəˈteriən/ (adj)}
{Enforcing strict obedience to authority.}
{He grew up in an authoritarian household with very strict rules.}
{Độc đoán.}

\vocabentry{Permissive}
{/pərˈmɪsɪv/ (adj)}
{Allowing excessive freedom of behavior.}
{Permissive parenting can sometimes lead to behavioral problems in children.}
{Dễ dãi/Buông thả.}

\vocabentry{Nurturing}
{/ˈnɜːrtʃərɪŋ/ (adj)}
{Caring for and encouraging development.}
{A nurturing environment is essential for a child's healthy growth.}
{Nuôi dưỡng/Chăm sóc.}

\vocabentry{Overprotective}
{/ˌoʊvərprəˈtektɪv/ (adj)}
{Tending to protect someone excessively.}
{Overprotective parents may prevent their children from developing independence.}
{Bao bọc quá mức.}

\vocabentry{Helicopter}
{/ˈhelɪkɑːptər/ (n)}
{A parent who is over-focused on their child.}
{Helicopter parenting can increase anxiety in college students.}
{Cha mẹ "trực thăng" - can thiệp quá sâu.}

\vocabentry{Rivalry}
{/ˈraɪvlri/ (n)}
{Competition for the same objective or superiority.}
{Sibling rivalry is quite common in large families.}
{Sự ganh đua.}

\vocabentry{Adoption}
{/əˈdɑːpʃn/ (n)}
{Legally taking another's child and bringing it up as one's own.}
{They are considering adoption as a way to grow their family.}
{Sự nhận con nuôi.}

\vocabentry{Foster}
{/ˈfɔːstər/ (adj)}
{Providing parental care to others' children for a period.}
{Foster parents provide a vital service for children in need.}
{Nhận nuôi tạm thời.}

\vocabentry{Infancy}
{/ˈɪnfənsi/ (n)}
{The state or period of early childhood.}
{Proper nutrition is especially critical during infancy.}
{Thời thơ ấu/Sơ sinh.}

\vocabentry{Toddler}
{/ˈtɑːdlər/ (n)}
{A young child who is just beginning to walk.}
{The toddler was excitedly exploring the backyard.}
{Trẻ chập chững biết đi.}

\vocabentry{Kindergarten}
{/ˈkɪndərɡɑːrtn/ (n)}
{School for children below the age of five.}
{Children learn through play and socialization in kindergarten.}
{Trường mẫu giáo.}

\vocabentry{Literacy}
{/ˈlɪtərəsi/ (n)}
{The ability to read and write.}
{Literacy rates have improved significantly over the last century.}
{Khả năng biết chữ.}

\vocabentry{Homeschooling}
{/ˌhoʊmˈskuːlɪŋ/ (n)}
{Education of children at home by their parents.}
{Homeschooling allows for a more personalized learning experience.}
{Tự học tại nhà.}

\vocabentry{Alternative}
{/ɔːlˈtɜːrnətɪv/ (adj)}
{Available as another possibility.}
{Many parents look for alternative education methods for their children.}
{Các lựa chọn thay thế.}

\vocabentry{Educational}
{/ˌedʒuˈkeɪʃənl/ (adj)}
{Relating to the provision of education.}
{Traveling can be a very educational experience for young people.}
{Có tính giáo dục.}

\vocabentry{Numeracy}
{/ˈnuːmərəsi/ (n)}
{The ability to understand and work with numbers.}
{Basic numeracy is an essential skill for everyday life.}
{Khả năng tính toán.}

\vocabentry{Scholarship}
{/ˈskɑːlərʃɪp/ (n)}
{A grant or payment made to support a student's education.}
{She was awarded a full scholarship to study at Oxford.}
{Học bổng.}

\vocabentry{Curriculum}
{/kəˈrɪkjələm/ (n)}
{The subjects comprising a course of study.}
{The national curriculum is being updated to include more technology.}
{Chương trình học.}

\vocabentry{Pedagogy}
{/ˈpedəɡɑːdʒi/ (n)}
{The method and practice of teaching.}
{Modern pedagogy emphasizes student-centered learning.}
{Nghiệp vụ sư phạm.}

\vocabentry{Standardized}
{/ˈstændərdaɪzd/ (adj)}
{Designed to be administered in a consistent manner.}
{Standardized tests are used to compare student performance across schools.}
{Chuẩn hóa.}

\vocabentry{Achievement}
{/əˈtʃiːvmənt/ (n)}
{A thing done successfully by effort or skill.}
{Winning the science fair was a great achievement for him.}
{Thành tích.}

\vocabentry{Vocational}
{/voʊˈkeɪʃənl/ (adj)}
{Relating to an occupation or employment.}
{Vocational training provides students with practical skills for work.}
{Thuộc về đào tạo nghề.}

\vocabentry{Internship}
{/ˈɪntɜːrnʃɪp/ (n)}
{Position of a student or trainee working to gain experience.}
{Doing an internship is a great way to network in your chosen field.}
{Kỳ thực tập.}

\vocabentry{Mentorship}
{/ˈmentɔːrʃɪp/ (n)}
{Period of time during which an experienced person helps another.}
{Mentorship can provide valuable guidance for young professionals.}
{Sự cố vấn.}

\vocabentry{Leadership}
{/ˈliːdərʃɪp/ (n)}
{The action of leading a group or organization.}
{Strong leadership is essential for any successful team.}
{Khả năng lãnh đạo.}

\vocabentry{Entrepreneur}
{/ˌɑːntrəprəˈnɜːr/ (n)}
{A person who organizes and operates a business.}
{He is a successful entrepreneur who started his own tech company.}
{Doanh nhân.}

\vocabentry{Financial}
{/faɪˈnænʃl/ (adj)}
{Relating to money or how it is managed.}
{Young people need to learn financial responsibility from an early age.}
{Thuộc về tài chính.}

\vocabentry{Independence}
{/ˌɪndɪˈpendəns/ (n)}
{The fact or state of being independent.}
{Moving away for university helps students develop their independence.}
{Sự tự lập.}

\vocabentry{Self-sufficient}
{/ˌself səˈfɪʃnt/ (adj)}
{Needing no outside help in satisfying basic needs.}
{Learning to cook is a step towards becoming self-sufficient.}
{Tự cung tự cấp/Tự lo.}

\vocabentry{Responsibility}
{/rɪˌspɑːnsəˈbɪləti/ (n)}
{The state or fact of having a duty.}
{Having a part-time job teaches students about responsibility.}
{Trách nhiệm.}

\vocabentry{Resilience}
{/rɪˈzɪliəns/ (n)}
{The capacity to recover quickly from difficulties.}
{Failure is an opportunity to build resilience and mental toughness.}
{Khả năng phục hồi.}

\vocabentry{Resourceful}
{/rɪˈsɔːrsfl/ (adj)}
{Ability to find quick ways to overcome difficulties.}
{A resourceful person can find a solution to almost any problem.}
{Tháo vát.}

\vocabentry{Ambition}
{/æmˈbɪʃn/ (n)}
{A strong desire to achieve something.}
{Her ambition is to become a world-renowned surgeon.}
{Tham vọng/Hoài bão.}

\vocabentry{Motivation}
{/ˌmoʊtɪˈveɪʃn/ (n)}
{Reason one has for acting in a particular way.}
{High grades are a strong motivation for many students.}
{Động lực.}

\vocabentry{Perseverance}
{/ˌpɜːrsəˈvɪərəns/ (n)}
{Persistence in doing something despite difficulty.}
{Perseverance is key to achieving long-term goals.}
{Sự kiên trì.}

\vocabentry{Dedication}
{/ˌdedɪˈkeɪʃn/ (n)}
{Quality of being committed to a task.}
{Success requires a lot of hard work and dedication.}
{Sự tận tụy.}

\vocabentry{Curiosity}
{/ˌkjʊriˈɑːsəti/ (n)}
{A strong desire to know or learn something.}
{Curiosity is the engine of scientific discovery.}
{Sự tò mò.}

\vocabentry{Creativity}
{/ˌkriːeɪˈtɪvəti/ (n)}
{The use of the imagination or original ideas.}
{Creativity should be encouraged in the classroom.}
{Sự sáng tạo.}

\vocabentry{Collaboration}
{/kəˌlæbəˈreɪʃn/ (n)}
{Working with someone to produce something.}
{The project was a successful collaboration between several departments.}
{Sự cộng tác.}

\vocabentry{Teamwork}
{/ˈtiːmwɜːrk/ (n)}
{Combined action of a group, especially when efficient.}
{Teamwork skills are essential for most professional roles.}
{Làm việc nhóm.}

\vocabentry{Networking}
{/ˈnetwɜːrkɪŋ/ (n)}
{Interacting with others to exchange information.}
{Networking events are a great way to find new job opportunities.}
{Việc kết nối quan hệ.}

\vocabentry{Advocacy}
{/ˈædvəkəsi/ (n)}
{Public support for a particular cause.}
{She is known for her advocacy of human rights.}
{Sự vận động/ủng hộ.}

\vocabentry{Empowerment}
{/ɪmˈpaʊərmənt/ (n)}
{Authority or power given to someone.}
{Education is a powerful tool for the empowerment of women.}
{Sự trao quyền.}

\vocabentry{Equality}
{/iˈkwɑːləti/ (n)}
{The state of being equal in status or rights.}
{We must strive for equality in the workplace.}
{Sự bình đẳng.}

\vocabentry{Equity}
{/ˈekwəti/ (n)}
{The quality of being fair and impartial.}
{Equity involves providing everyone with what they need to succeed.}
{Sự công bằng/Tính ngang bằng.}

\vocabentry{Humanity}
{/hjuːˈmænəti/ (n)}
{Human beings collectively.}
{We must show humanity and compassion to those in need.}
{Nhân tính/Nhân loại.}

\vocabentry{Compassion}
{/kəˈpæʃn/ (n)}
{Sympathetic pity for the sufferings of others.}
{He showed great compassion for the homeless people in his city.}
{Lòng trắc ẩn.}

\vocabentry{Altruism}
{/ˈæltruɪzəm/ (n)}
{Selfless concern for the well-being of others.}
{Her altruism led her to donate most of her savings to charity.}
{Lòng vị tha.}

\vocabentry{Integrity}
{/ɪnˈteɡrəti/ (n)}
{Quality of being honest and having strong morals.}
{He is a man of great integrity and always tells the truth.}
{Sự chính trực.}

\vocabentry{Honesty}
{/ˈɑːnəsti/ (n)}
{The quality of being honest.}
{Honesty is the best policy in any relationship.}
{Sự trung thực.}

\vocabentry{Authenticity}
{/ˌɔːθenˈtɪsəti/ (n)}
{The quality of being authentic.}
{Authenticity is key to building trust with your audience.}
{Tính xác thực.}

\vocabentry{Introspection}
{/ˌɪntrəˈspekʃn/ (n)}
{Examination of one's own mental processes.}
{Introspection can help you understand your motivations and desires.}
{Sự hướng nội/Soi xét nội tâm.}

\vocabentry{Perspective}
{/pərˈspektɪv/ (n)}
{Attitude toward or way of regarding something.}
{Traveling helps people gain a wider perspective on the world.}
{Góc nhìn.}

\vocabentry{Wisdom}
{/ˈwɪzdəm/ (n)}
{Quality of having experience and good judgment.}
{With age comes wisdom and life experience.}
{Sự thông thái.}

\vocabentry{Maturity}
{/məˈtʃʊrəti/ (n)}
{The state, fact, or period of being mature.}
{She showed great maturity in dealing with the difficult situation.}
{Sự trưởng thành.}

\vocabentry{Adaptability}
{/əˌdæptəˈbɪləti/ (n)}
{Quality of being able to adjust to new conditions.}
{Adaptability is an essential skill in today's fast-changing job market.}
{Khả năng thích nghi.}

\vocabentry{Agility}
{/əˈdʒɪləti/ (n)}
{Ability to move quickly and easily.}
{Mental agility is important for solving complex problems.}
{Sự nhanh nhẹn/linh hoạt.}

\vocabentry{Flexibility}
{/ˌfleksəˈbɪləti/ (n)}
{Quality of bending easily without breaking.}
{Remote work offers more flexibility for employees.}
{Sự linh động.}

\vocabentry{Passion}
{/ˈpæʃn/ (n)}
{Strong and barely controllable emotion.}
{He has a passion for music and spends most of his time playing the guitar.}
{Niềm đam mê.}

\vocabentry{Enthusiasm}
{/ɪnˈθuːziæzəm/ (n)}
{Intense and eager enjoyment or interest.}
{Her enthusiasm for her work is truly inspiring.}
{Sự nhiệt huyết.}

\vocabentry{Optimism}
{/ˈɑːptɪmɪzəm/ (n)}
{Hopefulness and confidence about the future.}
{Despite the challenges, she maintained her optimism for the project.}
{Sự lạc quan.}

\vocabentry{Hope}
{/hoʊp/ (n)}
{Feeling of expectation and desire for a certain thing.}
{There is always hope for a better future.}
{Hy vọng.}

\vocabentry{Believe}
{/bɪˈliːv/ (v)}
{Accept something as true.}
{You must believe in yourself to achieve your goals.}
{Tin tưởng.}

\vocabentry{Dream}
{/driːm/ (n)}
{Cherished aspiration or ideal.}
{Her dream is to travel around the world.}
{Ước mơ.}

\vocabentry{Aspiration}
{/ˌæspəˈreɪʃn/ (n)}
{Hope or ambition of achieving something.}
{Many young people have high aspirations for their careers.}
{Khát vọng.}

\vocabentry{Goal}
{/ɡoʊl/ (n)}
{Object of a person's ambition or effort.}
{Setting specific goals can help you stay motivated.}
{Mục tiêu.}

\vocabentry{Success}
{/səkˈses/ (n)}
{Accomplishment of an aim or purpose.}
{Success is the result of hard work and perseverance.}
{Sự thành công.}

\vocabentry{Failure}
{/ˈfeɪljər/ (n)}
{Lack of success.}
{Failure is a stepping stone to success if you learn from it.}
{Sự thất bại.}

\vocabentry{Experience}
{/ɪkˈspɪriəns/ (n)}
{Practical contact with and observation of facts.}
{You need several years of experience for this senior position.}
{Kinh nghiệm.}

\vocabentry{Knowledge}
{/ˈnɑːlɪdʒ/ (n)}
{Facts, information, and skills acquired.}
{Knowledge is power in today's information age.}
{Kiến thức.}

\vocabentry{Understanding}
{/ˌʌndərˈstændɪŋ/ (n)}
{Ability to understand something.}
{Mutual understanding is key to a healthy relationship.}
{Sự thấu hiểu.}

\vocabentry{Insight}
{/ˈaɪsaɪt/ (n)}
{Deep understanding of a person or thing.}
{The book provides valuable insight into the human mind.}
{Sự sáng suốt/nhìn sâu.}

